The two attention scenarios are easiest to read side by side. Making silence as salient as a rating changes what buyers think about only when they see no rating: the price of a displayed rating is unchanged and the price of silence falls, so the sellers who hid have no reason to. Renovation pays the same $b$ as before, so the threshold does not move. Full attention also changes what buyers think about when they do see a rating: a displayed good rating is worth $\rho\beta_1$ instead of $\rho\omega_1\beta_1$. The data identify only the product $\hat\rho=\rho\omega_1$, so they cannot say how much of a low $\hat\rho$ is few EPC-sensitive buyers (then nothing changes) and how much is inattention to displayed ratings (then the premium rises a great deal). The result is a bracket between the maintained reading and an upper bound set by the energy-cost anchor. Table~\ref{tab:att_scenarios} reports the scenarios for both segments.
